% s-088-p088-m001; original PDF 88
\Aphorism{s-088-p088-m001}{87}
Diejenigen, die scheinbar nichts auf das Urtheil Anderer geben, mögen nur bei sich suchen, in welchem Schlupfwinkel sich diese Abhängigkeit von Anderen verkrochen hat; dass sie noch zu finden ist, nehme ich auf mich. Eine sehr häufige Verkleidungsform ist die Consequenz, der „Charakter“, der das hundert-und-erste Mal nicht anders entscheiden zu dürfen glaubt, als er hundert Male entschieden hat; das Forum, das für diese Vergleichung des im Grunde Unvergleichbaren herangezogen wird, ist das Urtheil Anderer. Man hat aus sich selbst einen Präcedenzfall gemacht, den man Lügen zu strafen fürchtet, bei wem? bei aller Welt, bei den Anderen, denen man mit einer von Fall zu Fall wechselnden, individuellen Handlungsweise auffallen würde.
\par

